%% ma: L11-B07-0015
%% lop: 11
%% chuong: 2
%% bai: 7
%% dang: 11.7.7
%% loai: TLN
%% muc: TH
%% dap_an: "14,7"
%% nguon: "de-on-tap-giua-ki-1-toan-11-nam-2026-2027-truong-thpt-ngo-quyen-da-nang.pdf (M09, Đề ôn tập giữa kì 1 Toán 11 THPT Ngô Quyền, Đà Nẵng), Mã đề 0114, Phần III Câu 3"
%% kien_thuc: "Bài 7 (cấp số nhân)"
%% trang_thai: cho-duyet
%% ly_do: "Dạng mới đề xuất 11.7.7"
%% ghi_chu: "Hình con lắc ở gốc chỉ minh họa, không chứa dữ kiện toán học; bản kho không vẽ lại."
\begin{ex}
Một con lắc dao động tắt dần chậm với biên độ ban đầu là $20$ cm, biết cứ sau mỗi chu kì biên độ con lắc giảm đi một lượng bằng $5\%$ biên độ của lần dao động ngay trước đó. Hỏi sau $6$ lần dao động biên độ con lắc còn lại bao nhiêu centimét? (làm tròn kết quả đến hàng phần mười)
\shortans{14,7}
\loigiai{Sau mỗi chu kì biên độ nhân với $0{,}95$. Sau $6$ lần: $20\cdot0{,}95^6\approx14{,}7$ (cm).}
\end{ex}
